% !TEX program = xelatex
\documentclass{英语试卷}
\begin{document}
\examtitle{日照市 2022 级高三校际联合考试}{英语}
\par 日照市 2022 级高三校际联合考试\par
\par 英语试题卷\par
\begin{notice}
\par 1. 答卷前，考生务必将自己的姓名、准考证号填写在答题卡上。\par
\par 2. 回答选择题时，选出每小题答案后，用铅笔把答题卡上对应题目的答案标号涂黑。如需改动，用橡皮擦干净后，再选涂其他答案标号。回答非选择题时，将答案写在答题卡上，写在本试卷上无效。\par
\par 3. 考试结束后，将本试卷和答题卡一并交回。\par
\end{notice}
\parttitle{第一部分 听力(共两节，满分 30 分)}
\par 做题时，先将答案标在试卷上。录音内容结束后，你将有两分钟的时间将试卷上的答案转涂到答题卡上。\par
\sectitle{第一节(共 5 小题；每小题 1.5 分，满分 7.5 分)}
\par 听下面 5 段对话，每段对话后有一个小题，从题中所给的 A、B、C 三个选项中选出最佳选项。听完每段对话后，你都有 10 秒钟的时间来回答有关小题和阅读下一小题，每段对话仅读一遍。\par
\question{1. What will the man do this evening?}
\begin{choices}(3)
  \item Meet his co-worker.
  \item Do some washing.
  \item Attend a party.
\end{choices}
\question{2. Where can the man possibly find Peter?}
\begin{choices}(3)
  \item In the library.
  \item In the classroom.
  \item In the headmaster's office.
\end{choices}
\question{3. What is the man's attitude towards his problem?}
\begin{choices}(3)
  \item Sorry.
  \item Uncertain.
  \item Optimistic.
\end{choices}
\question{4. What do we know about the woman?}
\begin{choices}(3)
  \item She lost her car keys.
  \item She has poor eyesight.
  \item She forgot her license.
\end{choices}
\question{5. What time were Richard and Sarah due to arrive?}
\begin{choices}(3)
  \item At 7:30 p.m.
  \item At 8:30 p.m.
  \item At 9:30 p.m.
\end{choices}
\sectitle{第二节(共 15 小题;每小题 1.5 分,满分 22.5 分)}
\par 听下面 5 段对话或读白，每段对话或独白后有几个小题，从题中所给的 A、B、C 三个选项中选出最佳选项。听每段对话或独白前，你将有时间阅读各个小题，每小题 5 秒钟。听完后，各小题将给出 5 秒钟的作答时间。每段对话或独白读两遍。\par
\par 听第 6 段材料，回答第 6、7 题，\par
\question{6. What season might it be now?}
\begin{choices}(3)
  \item Winter.
  \item Autumn.
  \item Spring.
\end{choices}
\question{7. What gift will Linda receive?}
\begin{choices}(3)
  \item A mobile phone.
  \item A dictionary.
  \item A coat.
\end{choices}
\par 听第 7 段材料，回答第 8 至 10 题。\par
\question{8. Where does the conversation take place?}
\begin{choices}(3)
  \item In a gym.
  \item In a bank.
  \item In a shop.
\end{choices}
\question{9. What does the woman say about the place in the evenings?}
\begin{choices}(3)
  \item It's expensive.
  \item It's busy.
  \item It's dirty.
\end{choices}
\question{10. How much extra money will the man pay for an upgraded membership?}
\begin{choices}(3)
  \item £8.
  \item £10.
  \item £20.
\end{choices}
\par 听第 8 段材料，回答第 11 至 13 题。\par
\question{11. How long did Mike stay in Hong Kong?}
\begin{choices}(3)
  \item For about 5 days.
  \item For about 15 days.
  \item For about 20 days.
\end{choices}
\question{12. Why was May unable to get in touch with Mike?}
\begin{choices}(3)
  \item His phone was lost.
  \item His phone got broken.
  \item His phone number changed.
\end{choices}
\question{13. What did Mike do in Hong Kong?}
\begin{choices}(3)
  \item He went to the sea.
  \item He drove around with families.
  \item He worked part-time in Disneyland.
\end{choices}
\par 听第 9 段材料，回答第 14 至 17 题，\par
\question{14. What is the relationship between the speakers?}
\begin{choices}(2)
  \item Best friends. 3. Boss and secretary.
  \item Teacher and student.
\end{choices}
\question{15. Who is Katie Marsden?}
\begin{choices}(3)
  \item A friend of the man.
  \item A member of a school club.
  \item A champion cyclist.
\end{choices}
\question{16. What event will happen on Saturday?}
\begin{choices}(3)
  \item A meeting.
  \item A cycling race.
  \item An important exam.
\end{choices}
\question{17. What are the speakers mainly discussing?}
\begin{choices}(3)
  \item Doing homework.
  \item Joining a club.
  \item Playing sports.
\end{choices}
\par 听第 10 段材料，回答第 18 至 20 题，\par
\question{18. What was thrown into the sea?}
\begin{choices}(3)
  \item A cage.
  \item Meat.
  \item Fish.
\end{choices}
\question{19. Who is the speaker probably?}
\begin{choices}(3)
  \item A boat captain.
  \item A photographer.
  \item A tour guide.
\end{choices}
\question{20. Where is the speaker?}
\begin{choices}(3)
  \item In South Africa.
  \item In Australia.
  \item In America.
\end{choices}
\parttitle{第二部分 阅读理解（共两节，满分 50 分）}
\sectitle{第一节（共15小题；每小题2.5分，满分37.5分）}
\par 阅读下列短文，从每题所给的 A、B、C 和 D 四个选项中，选出最佳答案，\par
\passagelabel{A}
\par 21st ANNUAL TECH FOR GOOD HACKATHON\par
\par This year's hackathon(编程马拉松) introduces exciting new categories, offering you the chance to showcase your technical skills and creativity while addressing real-world challenges. Each category focuses on a specific social issue, so choose one (or more) and develop innovative solutions that make a difference!\par
\par DETAILS\par
\par\vspace{0.4em}
\begin{center}\begin{tabular}{{|l|p{\dimexpr\linewidth-3em-4\tabcolsep}|}}
\hline
Category 1:Health Tech Innovations & Focus on issues like:Developing apps or devices to improve healthcare access.Solutions for remote patient monitoring. \\
Category 2:Education Technology & Focus on issues like:Creating tools to enhance online learning.Bridging the digital divide in education. \\
Category 3: Sustainable Urban Solutions & Design a 60-120-second presentation videos showing technology-driven ideas to create smarter, greener cities. \\
\hline
\end{tabular}\end{center}
\vspace{0.4em}
\par RULES\par
\par ▶ Entries must be submitted by 11:59 PM May 15,2025.\par
\par ▶ Open to university students and young professionals, with prizes awarded to the top teams.\par
\par ▶ Projects can include software, hardware, or a combination of both. Presentation videos should be 60-120 seconds long.\par
\par Judging criteria include innovation, feasibility, and social impact. Entries containing copyrighted material will be disqualified.\par
\par Submit project documentation in PDF format; presentation videos in MP4, AVI, or similar formats.\par
\par HOW TO SUBMIT\par
\par Step 1: Registration Register your team on our official website and complete the PROJECT PROPOSALFORM\par
\par Step 2: Sign an agreement\par
\par Each team member must sign a Tech for Good participation agreement, allowing us to showcase your project. If you are under 18 years of age, a parent or guardian must complete the form on our website.\par
\par Step 3: Submission\par
\par Package your project files and presentation video into a single ZIP folder named "ProjectTitle\_TeamName" and upload it to the Tech for Good Submission Portal.\par
\question{21. What is the main purpose of the hackathon?}
\begin{choices}(2)
  \item To promote commercial software development.
  \item To teach basic programming skills to beginners.
  \item To provide job opportunities for young professionals.
  \item To encourage technological solutions to social problems.
\end{choices}
\question{22. Which of the following will lead to disqualification?}
\begin{choices}(2)
  \item Submitting a 100-second video.
  \item Using original music in the video.
  \item Uploading Word-format documentation.
  \item Combining hardware and software.
\end{choices}
\question{23. What must teams do before uploading their project files?}
\begin{choices}(2)
  \item Pay a registration fee online.
  \item Sign a participation agreement.
  \item Provide parents' signatures.
  \item Publish their proposal publicly.
\end{choices}
\passagelabel{B}
\par My friend had built a multimillion-dollar business from nothing, rising out of poverty. He cared for his aging parents, sent two siblings to college, and even helped them get married. Yet, despite these incredible achievements, his response was always the same whenever I praised him: "Not good enough. Not yet."\par
\par I know this mindset well, having grown up in an entrepreneurial family where my performance was constantly measured through comparisons. "Look at so-and-so's son and how well he's doing," my father would say. The thought of meeting these standards blinded me to the 90\% of things I was doing right. Instead, I became focused on the never-ending 10\% I hadn't achieved to prove myself, and if not, I would become bitter toward myself and those I was trying to match in performance. I would develop an inflated self-confidence that I would dominate all aspects of business or have no confidence in certain business activities.\par
\par When I didn't feel competitive enough, I lost opportunities and friendships when there was a better approach for collaboration and a win-win. I often didn't know why I was competing; I just wanted to prove to my father that I was better than the other guy, and he was wrong. Later in my life, I came to realize that life can be competitive, but comparison is not the right approach to growth. It's the enemy of self-acceptance, breeding feelings of inadequacy. When we compare ourselves, we might create slightly better versions than others, but it comes at the expense of our authenticity.\par
\par When you accept yourself, you'll realize that you only need to meet your natural and basic expectations, not anyone else's. That's when you begin to live authentically and give others the space to grow into their best selves within their abilities and circumstances. An oak tree doesn't strive to become a smaller banana tree or a much larger mountain. It grows into its most natural and beautiful form, providing shade and strength for those who need it.\par
\question{24. What does the author want to show by telling his friend's story?}
\begin{choices}(2)
  \item Family support is the key to success.
  \item Poverty motivates greater achievements.
  \item Financial success doesn't mean a happy life.
  \item Achievement doesn't guarantee self-acceptance.
\end{choices}
\question{25. What does the underlined word "inflated" in paragraph 2 mean?}
\begin{choices}(2)
  \item Overstated.
  \item Humble
  \item Stable
  \item Unsurprising.
\end{choices}
\question{26. What is a consequence of comparison according to the author?}
\begin{choices}(2)
  \item Motivation decline.
  \item Self-worth decrease.
  \item Personality disorder.
  \item Superiority complex.
\end{choices}
\question{27. What does the last paragraph imply?}
\begin{choices}(2)
  \item Constant comparison fuels one's ambition.
  \item Competition contributes to success.
  \item Self-acceptance enables authentic growth.
  \item Social acceptance determines one's value.
\end{choices}
\passagelabel{C}
\par Seed libraries emerged as a revolutionary concept nearly a generation ago. The earliest example is BASIL, the Bay Area Seed Interchange Library. In 2000, University of California Berkeley closed its campus farm to make room for research into genetically modified corn. This left the farm's collection of seeds homeless. Activists persuaded the Berkeley Ecology Center to sponsor a lending library that would make the seeds available to local gardeners, who in turn would harvest and return seeds from the crops they grew to make possible the next year's distribution.\par
\par From there, the seed library concept took off, spreading throughout all 50 states. Not only do seed libraries offer customers a free start in gardening, they can also make available seeds adapted to local growing conditions. Garden retailers(零售商) tend to stock nationally distributed lines of best-selling seeds that are chosen for their ability to grow over a broad area. A seed library, which isn't driven by the need for profit, can focus on what performs best in the local soils and climate.\par
\par Besides distributing seeds, seed libraries also encourage seed saving among their customers, so that the loans can be repaid at the end of the growing season. This is why the vegetable, fruit and flower varieties that seed libraries distribute are typically open pollinated(授粉). That is, when pollinated naturally by insects, birds, or wind, the plants bear seeds that produce a new generation closely resembling the parents. Hybrids, the mainstay of conventional seed producers, do not "come true" in this fashion, and so they're avoided by most seed libraries.\par
\par This has a couple of significant impacts. First, seed libraries tend to emphasize old-fashioned, heirloom(原种) plant varieties that originated before the modern emphasis on hybrid fruits and vegetables. In this way, seed libraries play a significant part in helping to keep such heirlooms available to gardeners. Then, because gardeners typically save seed from the stronger and more fruitful samples, they tend to preserve and reproduce the best adapted individuals. This practice encourages the gradual evolution of an ever-more locally adapted species. In this way, seed libraries are becoming ever more relevant to the needs of gardeners. 微信公众号 IAI English\par
\question{28. What distinguishes seed libraries from garden retailers?}
\begin{choices}(2)
  \item They focus on locally adapted seeds.
  \item They offer seeds at a lower price.
  \item They prioritize profit over sustainability.
  \item They distribute best-selling seeds.
\end{choices}
\question{29. What is stressed about open-pollinated plants in paragraph 3?}
\begin{choices}(2)
  \item Their high output.
  \item Their distinct appearance.
  \item Their rich varieties.
  \item Their biological advantages.
\end{choices}
\question{30. What is one important role of seed libraries?}
\begin{choices}(2)
  \item Promoting the study of hybrid fruits and vegetables.
  \item Ensuring the availability of heirloom plant varieties.
  \item Making efforts to develop entirely new plant species.
  \item Encouraging gardeners to change traditional methods.
\end{choices}
\question{31. Which of the following can be the best title for the text?}
\begin{choices}(2)
  \item Gardeners Starve for Best Seeds
  \item The Rise and Fall of Seed Libraries
  \item Seed Libraries Help Gardening Grow
  \item The Future of Hybrid-plant Gardening
\end{choices}
\passagelabel{D}
\par Language models like Claude aren't programmed directly by humans — instead, they're trained on large amounts of data. During that training process, they learn their own strategies to solve problems. These strategies are encoded in the billions of computations a model performs for every word it writes. They arrive incomprehensible to us, the model's developers. Put simply, we have limited insight into how models do most of the things they do.\par
\par Knowing how models like Claude think would allow us to have a better understanding of their abilities, as well as help us ensure that they're doing what we intend them to. We take inspiration from the field of neuroscience(神经科学), which has long studied the messy insides of thinking organisms, and try to build a kind of AI microscope that will let us identify patterns of activity and flows of information. There are limits to what you can learn just by talking to an AI model — after all, humans (even neuroscientists) don't know all the details of how our own brains work. So we look inside.\par
\par The researchers conduct two new studies that represent progress on the development of the "microscope", and the application of it to see new "AI biology". In the first study, we extend our prior work locating interpretable concepts ("features") inside a model to link those concepts together into computational "circuits"(电路), revealing parts of the pathway that transforms the words that go into Claude into the words that come out. In the second, we look inside Claude 3.5 Haiku, performing deep studies of simple tasks representative of ten crucial model behaviors.\par
\par The studies shed light on a part of what happens when Claude responds to these prompts(提示). Claude sometimes thinks in a conceptual space that is shared between languages, suggesting it has a kind of universal "language of thought." Claude will plan what it will say many words ahead, and write to get to that destination. This is powerful evidence that even though models are trained to output one word at a time, they may think on much longer horizons to do so.\par
\question{32. What do we know about Claude from paragraph 1?}
\begin{choices}(2)
  \item Claude's outputs are predictable.
  \item Claude learns autonomously from data.
  \item Claude can collect large amounts of data.
  \item Claude's decisions are easily explainable.
\end{choices}
\question{33. Why does the author mention neuroscience in paragraph 2?}
\begin{choices}(2)
  \item To explore how AI models work.
  \item To show the limits of AI models.
  \item To compare AI with human brains.
  \item To study how AI is programmed.
\end{choices}
\question{34. What can we learn about the studies?}
\begin{choices}(2)
  \item They measure AI's word-output speed.
  \item They record AI's rapid development.
  \item They reveal AI's inner computations.
  \item They show AI's practical application.
\end{choices}
\question{35. What does the last paragraph mainly talk about?}
\begin{choices}(2)
  \item Claude's unique thinking mode.
  \item Claude's training data source.
  \item Claude's word processing efficiency.
  \item Claude's language learning ability
\end{choices}
\sectitle{第二节（共5小题；每小题2分，满分12.5分）}
\par 根据短文内容，从短文后的选项中选出能填入空白处的最佳选项。选项中有两项为多余选项，\par
\par In today's fast-paced society, the importance of time management is undeniable. Everyone strives to learn and work more efficiently within the limited time they have. \blanknum{36} Here're some tips!\par
\par When you're busy, every moment in the day is precious. \blanknum{37} When you're taking the bus or you're waiting for a lecture, you can get something done. Do some research for an important project and pin some online resources. Read part of the studying material you need to cover for the tests.\par
\par Make to-do lists every single day! Write a monthly plan and include all tests, papers and courses you need to cover. Then, translate that plan into weekly goals and create daily to-do lists that help you achieve them. \blanknum{38} Start with the most important one. Use different colors to indicate the importance of the tasks on your list and act accordingly. If you have to postpone some of the tasks, you'll be relieved you're doing that with the less important ones.\par
\par If you find yourself struggling with the amount of work you have, talk to your professors immediately. American professors have office hours when students can come ask questions and get help. They want you to succeed, so ask for help! There are also many campus resources that you can use, such as tutors and the writing center. \blanknum{39}\par
\par Learn how to say no. You always enjoy hanging out with your friends, but sometimes that's not possible. \blanknum{40} It's okay to say no. You have your life in front of you; you can still hang out with them. At this moment, it's important to do what's right.\par
\begin{choices}(1)
  \item Have your tablet with you at all times.
  \item Have too many tasks on your to-do list?
  \item Never hesitate to get help when needed.
  \item But how can this be achieved in universities?
  \item But why should we care about efficiency at all?
  \item Maybe you have a paper to write or you're simply too tired.
  \item With your own effort, you will find tasks are easier to finish.
\end{choices}
\parttitle{第三部分 语言知识运用（共两节，满分 30 分）}
\sectitle{第一节（共15小题；每小题1分，满分15分）}
\par 阅读下面短文，从短文后各题所给的 A、B、C 和 D 四个选项中，选出可以填入空白处的最佳选项，\par
\par A United Airlines pilot has gained widespread attention online after personally overseeing the order and distribution of 30 pizzas to 155 passengers on his flight. Now nicknamed the "Pizza Pilot," Scott Wardle, who has flown for three decades, \blanknum{41} to ensure passengers felt cared for during what could have been a(n) \blanknum{42} experience.\par
\par The \blanknum{43} unfolded during a flight from Boise to Phoenix via Houston. Approximately 45 minutes into the \blanknum{44} a passenger fell unconscious, \blanknum{45} Wardle to divert(改变) the flight to Albuquerque for medical \blanknum{46} He landed the plane around 11 p. m., just as most airport dining options had closed.\par
\par After that, Wardle learned the flight would face additional \blanknum{47} because the cabin crew had reached their maximum working hours and would need to be \blanknum{48}. He began to seek a way to \blanknum{49} his stranded(滞留) passengers. He decided on pizza, \blanknum{50} it was a universally liked option. He \blanknum{51} ordered dozens of pizzas with his phone and began serving them with plates and napkins. The heartwarming \blanknum{52} quickly went viral, with photos and videos capturing him in his pilot uniform \blanknum{53} slices.\par
\par United Airlines \blanknum{54} his actions and offered to reimburse(补偿) him for the pizzas while praising him as a shining example of their \blanknum{55} to customer care. "He waited until every passenger received their slice before enjoying a pizza himself," a spokesperson noted.\par
\cz{41}{stepped up}{happened}{went out}{pretended}
\cz{42}{exhausting}{unique}{embarrassing}{unpleasant}
\cz{43}{adventure}{encounter}{incident}{conflict}
\cz{44}{search}{exploration}{journey}{experiment}
\cz{45}{allowing}{forcing}{reminding}{enabling}
\cz{46}{equipment}{research}{insurance}{assistance}
\cz{47}{delays}{challenges}{breakdowns}{expenses}
\cz{48}{dismissed}{replaced}{controlled}{examined}
\cz{49}{gather}{persuade}{transport}{comfort}
\cz{50}{admitting}{reasoning}{confirming}{announcing}
\cz{51}{definitely}{immediately}{eventually}{apparently}
\cz{52}{ceremony}{reunion}{scene}{passenger}
\cz{53}{handing out}{wrapping up}{putting away}{eating up}
\cz{54}{recognized}{detected}{evaluated}{described}
\cz{55}{response}{approach}{attitude}{commitment}
\sectitle{第二节（共10小题；每小题1.5分，满分15分）}
\par 阅读下面短文，在空白处填入 1 个适当的单词或括号内单词的正确形式。\par
\par Gulbar Herem enjoyed sitting by the window while taking the train back to her home in Qira county, Hotan prefecture. \blanknum{56} the window, sand dunes(沙丘) of the Taklimakan Desert in the Xinjiang Uygur autonomous region flashed by against the blue sky. "We are actually traveling through the desert \blanknum{57} is literally called the place with no return." Gulbar said.\par
\par The Hotan-Ruoqiang railway, \blanknum{58} key national project, has made it much more convenient to travel to and from southern Xinjiang. The price is also fair, with the cheapest train ticket \blanknum{59} (sell) at about 170 yuan.\par
\par With the opening of the railway comes more opportunities \blanknum{60} (see) other places in the world. "Now, people in Lop can easily travel elsewhere by train to find new opportunities \blanknum{61} just go sightseeing in other places. In fact, the new line isn't only about adding a way of traveling, it has made people see China's \blanknum{62} (determine) to develop Xinjiang," said Akbar Turson, a passenger from Lop county.\par
\par According to China Railway Urumqi Group, one year after the Hotan-Ruoqiang railway \blanknum{63} (put) into use, more than 360,900 passengers have traveled on it. Besides making traveling \blanknum{64} (easy), the Hotan-Ruoqiang railway has also improved logistic(物流)service and put local economic development on the fast track. What's more, by linking the railway between Ruoqiang and Golmud, Qinghai province, the railway has also created another logistic route \blanknum{65} (connect) southern Xinjiang with other parts of China.\par
\parttitle{第四部分 写作（共两节，满分 40 分）}
\sectitle{第一节（满分15分）}
\par 假定你是某国际学校的学生李华，你校将对图书馆进行升级改造，面向全校师生征求意见。请你给校长Mr. David写一封信，内容包括：\par
\par (1) 图书馆存在的问题;\par
\par (2) 你的建议。 注意：\par
\par (1) 写作词数应为 80 个左右;\par
\par (2) 请按如下格式在答题卡的相应位置作答。\par
\par Dear Mr. David,\par
\sectitle{第二节（满分25分）}
\par 阅读下面材料，根据其内容和所给段落开头语续写两段，使之构成一篇完整的短文。\par
\par Being a single mother had been difficult for Kate ever since her husband passed away. Balancing a job that barely paid enough while taking care of her daughter, Holly had been a daily struggle. Money was always tight, and there were many sacrifices, but Kate was determined to make sure Holly never felt like she was missing out on love or happiness. She had been saving up for this dinner for months, wanting to give Holly a wonderful birthday.\par
\par As they stepped into the restaurant, Holly's excitement was nearly impossible to contain. Her small hands flipped through the menu. Before she could make a choice, a group of waiters appeared, singing "Happy Birthday" as they carried a towering cake. Holly's eyes grew wide with wonder, her smile stretching from ear to ear. Seeing her daughter's face light up made every penny Kate had saved worth it. For a moment, Kate forgot about the financial struggles, the hard work, and the sacrifices. All that mattered was Holly's happiness.\par
\par As Holly tasted her cake, her face glowing with delight, a waiter approached Kate quietly and gestured for her to step aside. "Ma'am, the cost is going to double," the waiter said. "Double? Why?" she responded, panic bubbling up in her voice. "The recipe for the cake has changed, and you will have to pay 100 now," the waiter explained. Her throat tightened. "100 for a cake?! But . I can't afford it," she whispered. "I'm sorry, ma'am. But you'll have to pay," the waiter said.\par
\par With trembling hands, Kate returned to the table, forcing a smile for Holly. When it was time to leave, Kate felt her heart pounding as she asked Holly to wait by the door. Her mind raced, trying to figure out what she could do. Her hands were trembling as she spoke. But to her surprise, the waiter responded, "The bill has already been paid by a woman." Tears of relief sprang to Kate's eyes. "Do you know who she was?" Kate asked, her voice thick with emotion.\par
\par (1) 续写词数应为 150 左右;\par
\par (2) 请按如下格式在答题卡的相应位置作答，\par
\writingpara{Paragraph 1: The waiter pointed toward the exit.}
\writinglines{8}
\writingpara{Paragraph 2: As the woman walked away, Kate stood there, lost in thought.}
\writinglines{8}
\end{document}
